%% ============================================================================
%% Fichier  : annexes/Ann_C_templates.tex
%% Titre    : Modèles opérationnels
%% Version  : 1.0 -- novembre 2026
%% ============================================================================
%% Alignée sur les annexes des autres cours de l'auteur (ACSI, INF3226).
%% ============================================================================

\chapter{Modèles Opérationnels}
\label{ann:templates}

Cette annexe rassemble les modèles à reproduire en mission. Chacun renvoie
au chapitre où il est expliqué.

\section*{Fiche de chaîne de custody}

\begin{tableaucours}{L{4.0cm}L{8.4cm}}{%
    \tableauHeader{Rubrique} &
    \tableauHeader{Contenu}
}
Identification & Numéro de pièce, description, numéro de série, état apparent \\
Saisie & Lieu, date et heure, personne, cadre légal (réquisition, perquisition) \\
Acquisition & Outil et version, format, empreintes SHA-256 de l'original et de l'image \\
Scellé & Numéro, date, signataires \\
Mouvements & Date et heure, remis par, remis à, motif, empreinte vérifiée \\
Restitution ou destruction & Décision, date, signataires \\
\end{tableaucours}

\section*{Procès-verbal d'acquisition}

\begin{itemize}
    \item Cadre~: autorité, réquisition, mission, personnes présentes.
    \item Support~: description, photographies, état allumé ou éteint.
    \item Ordre de volatilité respecté (RFC~3227) et justification des choix.
    \item Outils, versions, paramètres~; empreintes calculées et vérifiées.
    \item Incidents et écarts à la procédure, consignés au moment où ils surviennent.
\end{itemize}

\section*{Plan du rapport forensique}

\begin{tableaucours}{L{3.6cm}L{8.8cm}}{%
    \tableauHeader{Section} &
    \tableauHeader{Contenu}
}
Synthèse & Deux à quatre pages, sans jargon~: questions, réponses, niveau de confiance \\
Mission & Contexte et questions posées, telles quelles \\
Méthode & Normes suivies, outils et versions, ordre des opérations \\
Constatations & Faits numérotés, horodatés, avec leur source et leur annexe \\
Analyse & Chronologie, niveau de confiance, hypothèses alternatives, limites \\
Conclusions & Réponses aux questions~; la qualification juridique relève de l'autorité judiciaire \\
Annexes & Pièces, empreintes, formulaires de custody \\
\end{tableaucours}

\section*{Relecture avant dépôt}

\begin{itemize}
    \item Chaque phrase des constatations est-elle vérifiable telle quelle par un autre expert~?
    \item Chaque analyse a-t-elle un niveau de confiance et une hypothèse alternative examinée~?
    \item Chaque pièce citée a-t-elle son empreinte en annexe~?
    \item Le texte évite-t-il toute qualification pénale et toute attribution non établie~?
\end{itemize}
Voir Ch.~\ref{chap:rapport}.
